%==============================================================================
% Section 7: Physical Interpretation and Outlook
%==============================================================================
\section{Physical Interpretation and Outlook}
\label{sec:discussion}

This section clarifies what has been completed and what has not yet been
settled.

\subsection{What Has Been Completed in homogeneous bulk geometric sector}

The principal items established in this paper at the classical quadratic level
of the homogeneous minisuperspace geometric sector are summarised in
\S\ref{sec:conclusion}.  Here we add a remark on the scope limitation.

``Quadratic block organisation'' means that the homogeneous bulk sector can be
decomposed into spin-2, physical spin-1, auxiliary spin-1, and scalar/torsion
subsectors.  Note, however, that the scalar/torsion subsector is not fully
diagonal in the $(r,\eta,V)$ variables: at quadratic order the radial mode mixes
with axial torsion, while the $V$ direction remains separate at the isotropic
point.

It is important to emphasise that \emph{inhomogeneous modes are not addressed
here}.  This paper concerns the classical EFT structure of the homogeneous bulk
geometric sector.

\subsection{What Is Not Claimed Here}

The following claims are deliberately withheld.
\begin{itemize}
  \item Full stability including inhomogeneous modes beyond minisuperspace.
  \item Identification of the generalised eigenvalue $\lambda_\phys$ with a
        universal mass squared.
  \item A claim that the full 4-dimensional deformed geometry constitutes a new
        Lie group manifold.
  \item Complete action generation including arbitrary non-invariant polynomials.
\end{itemize}
These are not in contradiction with the results of this paper; they are
consciously outside the stated scope.

\subsection{Relation to the Thermal Interpretation of the $S^1$ Direction}

Since paper~I \cite{paper1}, the $\Sone$ direction has been regarded as a natural
connection point to a thermal interpretation.  Although this paper does not
explicitly discuss that connection, the final-basis organisation of the homogeneous
geometric sector provides a firmer foundation for future discussions of:
\begin{itemize}
  \item the stable sector of the background geometry;
  \item control parameters coupled to thermal deformations;
  \item the separation of bulk geometric response from thermally dependent
        responses.
\end{itemize}

\subsection{Toward Future Extensions}

The future tasks for this model are the following.
\begin{enumerate}
  \item Introduction of inhomogeneous fluctuations including local degrees of
        freedom.
  \item Settings involving boundaries and interfaces.
  \item Connection between homogeneous bulk modes and localised modes.
  \item Bridges to physical observables including thermal and Lorentzian
        interpretations.
\end{enumerate}
